\section{Introduction}
\label{sec:intro}

Text-to-SQL systems translate natural-language questions into SQL queries,
enabling non-expert users to query relational databases directly.
Enterprise analytics platforms, business intelligence tools, and conversational
data assistants all represent deployment contexts for such systems.
Accurate, low-latency SQL generation is therefore a practical prerequisite
for wide adoption.

Recent state-of-the-art systems rely on large proprietary language models
such as GPT-4~\citep{openai2023gpt4}, achieving execution accuracy (EX) of
up to 0.866 on the Spider~1.0 benchmark~\citep{gao2023dailsql}.
These results, however, come with significant barriers: cloud API costs,
data privacy constraints, network latency, and the requirement of continuous
internet connectivity.
A growing class of applications---clinical decision support, on-device
analytics, air-gapped enterprise environments---cannot use cloud-hosted models.

Despite this progress, no rigorous study of a fully local, training-free
pipeline has examined the interaction between retrieval encoder quantization
and downstream SQL generation quality.
Prior quantization research focuses almost exclusively on the language model
itself~\citep{dettmers2022llmint8,xiao2023smoothquant,
shen2020qbert,zafrir2019q8bert};
the schema-linking retrieval encoder---a critical component of
retrieval-augmented Text-to-SQL~\citep{gao2023dailsql,eben2025rasl}---has
not been examined.
This gap is practically important: an encoder that degrades under
quantization would propagate retrieval errors directly into SQL generation
accuracy, invalidating the benefit of model compression.

This paper addresses that gap with four contributions.
First, it presents a systematic ablation of five pipeline components on
Spider~1.0 (1{,}034 development questions), establishing statistically
significant ($p < 0.001$, McNemar's test) accuracy improvements for each.
Second, it provides the first INT8 quantization study for a schema-retrieval
encoder in a Text-to-SQL pipeline, demonstrating 78.2\% RAM reduction and
15.2\% latency improvement with zero retrieval-quality degradation, verified
through an 84-experiment per-layer and per-sublayer sensitivity analysis.
Third, validation on Spider~2.0-Lite~\citep{lei2025spider2}---a substantially
harder enterprise benchmark---confirms that both the pipeline design and the
quantization safety finding generalize beyond Spider~1.0.
Fourth, the complete, reproducible implementation runs locally on consumer
hardware without cloud API access or model fine-tuning.

Figure~\ref{fig:overview} provides a high-level view of the three
contributions in a single diagram.
Panel~(a) shows the end-to-end pipeline: a natural-language question is
encoded by \texttt{arctic-embed-m}, matched against a pre-computed
\texttt{sqlite-vec} index to retrieve the top-5 relevant tables, and
fed---together with sql2skeleton few-shot examples---to a locally running
Qwen3.5~9B LLM that generates the final SQL under self-consistency with
two-pass error correction, achieving EX~=~0.824 on Spider~1.0.
Panels~(b) and~(c) contrast the FP32 baseline encoder (418\,MB RAM,
64.1\,ms/query) with the INT8-quantized variant (91\,MB RAM, 54.4\,ms/query):
the 84-experiment per-block sensitivity analysis confirms $\Delta$R@5\,=\,0
and $\Delta$MRR\,=\,0, establishing INT8 as the Pareto-optimal profile.

\begin{figure*}[!t]
  \centering
  \includegraphics[width=\linewidth]{fig0_overview.pdf}
  \caption[System overview: pipeline, FP32 baseline, and INT8-quantized encoder]{
    Overview of the proposed system.
    \textbf{(a)}~End-to-end Text-to-SQL pipeline: a natural-language question
    is encoded by the arctic-embed-m retrieval encoder, matched via cosine ANN
    search against a pre-computed \texttt{sqlite-vec} index to retrieve the
    top-5 relevant tables and their DDL, then combined with
    sql2skeleton-selected few-shot examples into a DAIL-SQL prompt for
    Qwen3.5~9B running locally under self-consistency ($N$\,=\,3) with
    two-pass error correction.
    The full pipeline achieves EX\,=\,0.824 on Spider~1.0 and EX\,=\,0.111
    on Spider~2.0-Lite.
    \textbf{(b)}~FP32 encoder baseline: all 12 encoder blocks and the pooler
    run in FP32 (418\,MB RAM, 64.1\,ms encoding latency,
    R@5\,=\,0.9952, MRR\,=\,0.9513).
    The sensitivity heat strip at the bottom shows the per-block
    $\Delta$R@5 for all 84 INT8 experiments.
    \textbf{(c)}~INT8-quantized encoder (this work): all \texttt{nn.Linear}
    layers are quantized to INT8 via \texttt{torch.quantization.quantize\_dynamic}
    (token embeddings remain FP32 as \texttt{nn.Embedding} is not targeted).
    RAM drops to 91\,MB ($-$78.2\%), latency to 54.4\,ms ($-$15.2\%), and
    $\Delta$R@5\,=\,0 with $\Delta$MRR\,=\,0---identical retrieval quality
    to FP32 at one-quarter the memory footprint.
    The runtime quantization engine is \texttt{qnnpack} on ARM64
    (Apple~M-series) and \texttt{fbgemm} on x86-64, selected automatically.}
  \label{fig:overview}
\end{figure*}

The remainder of this paper is organized as follows.
Section~\ref{sec:background} reviews related work.
Section~\ref{sec:pipeline} describes the pipeline architecture.
Section~\ref{sec:quantization} presents the quantization analysis.
Section~\ref{sec:experiments} reports ablation and comparison experiments.
Section~\ref{sec:discussion} discusses limitations and design implications.
Section~\ref{sec:conclusion} concludes.
